\documentclass[11pt]{article}
\usepackage[margin=1in]{geometry}
\usepackage{amsmath}
\usepackage{booktabs}
\usepackage{tikz}
\usepackage{xcolor}
\usepackage{listings}
\usepackage{hyperref}

% --- code listing style for Verilog-ish path snippets ---
\lstset{
  basicstyle=\ttfamily\small,
  breaklines=true,
  frame=single,
  columns=fullflexible,
  keepspaces=true,
}

\title{PPU Timing Closure: Path-1 Pipelining}
\author{}
\date{\today}

\begin{document}
\maketitle

\section{Original Architecture (Before Fix)}

The PPU consumes one $16\times1$ MAC result row per call and performs a
FlashAttention-style streaming softmax over $64$ tiles. Path-1 builds the
running row maximum and exponent sum; Path-2 replays the tiles to emit the
quantized softmax output. Each row is processed over two half-cycles (low/high)
so that a single folded engine of $8$ \texttt{LPW\_expo} units is time-shared
across the $16$ lanes (initiation interval $\text{II}=2$ cycles/row).

The problem: Path-1 executed its \emph{entire} datapath in a single clock cycle:

\begin{lstlisting}
vec_max  ->  new_max compare  ->  LPW_expo(delta)  ->  x prev_sum  ->  + exp_sum  ->  sum_buffer
\end{lstlisting}

This chains a max-reduction, an exponent unit, a multiplier and two adders in
series within one cycle.

\section{First Synthesis Result}

Synthesis at a $10$~ns clock period ($100$~MHz) reported a setup violation on
this Path-1 cone:

\begin{center}
\begin{tabular}{lr}
\toprule
Quantity & Value (ns) \\
\midrule
Data required time & $10.82$ \\
Data arrival time  & $15.14$ \\
\midrule
\textbf{Slack (VIOLATED)} & $\mathbf{-4.32}$ \\
\bottomrule
\end{tabular}
\end{center}

The arrival time of $\sim15$~ns confirms $\sim14$~ns of combinational logic
(after subtracting the $1$~ns ideal clock-network delay) in a single cycle.
Breaking the reported path down by sub-block:

\begin{center}
\begin{tabular}{llr}
\toprule
Path window (ns) & Block & $\Delta$ (ns) \\
\midrule
$1.66 \to 5.31$  & \texttt{vec\_max} (row-max reduction) & $3.65$ \\
$5.31 \to 7.15$  & \texttt{new\_max} compare + mux glue   & $1.84$ \\
$7.15 \to 10.91$ & \texttt{LPW\_expo} (sub + mult + add)  & $3.76$ \\
$11.57 \to 13.64$& \texttt{mult\_132} ($\text{prev\_sum}\times\delta$) & $2.07$ \\
$13.92 \to 14.64$& \texttt{add\_166} (accumulate)         & $0.72$ \\
\bottomrule
\end{tabular}
\end{center}

The single largest block is the row-max reduction (\texttt{vec\_max},
$3.65$~ns), but note that \texttt{vec\_max} and \texttt{LPW\_expo} together are
already $\sim10.6$~ns, so no single register cut can close timing---at least two
cuts are required.

\section{Fix: 4-Stage Path-1 Pipeline}

Path-1 was restructured into a four-stage pipeline. Throughput is unchanged
(still $\text{II}=2$ cycles/row); only the latency increases by $2$ cycles.

\begin{center}
\begin{tabular}{cllr}
\toprule
Stage & Name & Work & Est.\ logic (ns) \\
\midrule
M  & Max      & \texttt{vec\_max} + \texttt{new\_max}; read max/sum buffers & $\sim4.5$ \\
E0 & Exp-low  & folded \texttt{LPW} low $8$ lanes + delta-exp factor       & $\sim5.0$ \\
E1 & Exp-high & folded \texttt{LPW} high $8$ lanes + current exp sum        & $\sim5.0$ \\
A  & Accum    & $\text{prev\_sum}\times\delta + \text{exp\_sum}$; writeback & $\sim3.5$ \\
\bottomrule
\end{tabular}
\end{center}

\noindent
(The per-stage values above are \emph{estimates}; the synthesizer re-optimizes
across the new register boundaries, so the true per-stage delays will come from
the re-run reported in Section~\ref{sec:after}.)

\paragraph{Pipeline overlap.} With $\text{II}=2$, the $16$ rows of a tile launch
on cycles $0,2,\dots,30$ and each writes back $3$ cycles after launch. At most
$2$ rows are ever in flight, so only two row buffers are needed. The extra
latency appears \emph{once per tile} (the last row writes back $2$ cycles later)
and is absorbed by the existing inter-tile gap (MAC runs at $35$~cycles/tile vs.\
PPU's $32$-cycle drain).

\begin{center}
\begin{tikzpicture}[font=\ttfamily\footnotesize, x=0.85cm, y=0.6cm]
  % axis
  \draw[->] (0,0) -- (14,0) node[right]{cycle};
  \foreach \c/\lab in {0/0,2/2,4/4,6/6,8/8,13/33} {
    \draw (\c,0.1) -- (\c,-0.1) node[below]{\lab};
  }
  % row0: M E0 E1 A on cycles 0,1,2,3
  \node[anchor=west] at (-1.8,1) {row0:};
  \node at (0,1){M}; \node at (1,1){E0}; \node at (2,1){E1}; \node at (3,1){A};
  % row1: launched at cycle 2 (II=2)
  \node[anchor=west] at (-1.8,2) {row1:};
  \node at (2,2){M}; \node at (3,2){E0}; \node at (4,2){E1}; \node at (5,2){A};
  % row2
  \node[anchor=west] at (-1.8,3) {row2:};
  \node at (4,3){M}; \node at (5,3){E0}; \node at (6,3){E1}; \node at (7,3){A};
  % dots
  \node at (9,3.4){$\cdots$};
  % row15: launched at cycle 10
  \node[anchor=west] at (-1.8,4) {row15:};
  \node at (10,4){M}; \node at (11,4){E0}; \node at (12,4){E1}; \node at (13,4){A};
\end{tikzpicture}
\end{center}

\paragraph{Key design choices.}
\begin{itemize}
  \item \textbf{Kept the 8-LPW folding.} Un-folding to $16$ \texttt{LPW\_expo}
        units would remove the low/high split but roughly triple the exponent
        multiplier area; the folded engine is the better area$\times$time point.
  \item \textbf{Folded-engine half-select driven by the pipeline's own valid
        bits} (\texttt{e1\_valid}), not the live \texttt{valid} input, so the
        high-half selection stays correct when the feeder pauses at a tile's
        last row.
  \item \textbf{Removed the PPU-internal \texttt{row\_buffer}.} The MAC tile's
        \texttt{latch\_array} already stores the whole tile, so the top-level
        feeder now holds \texttt{mac\_result\_vec} stable across both half-cycles
        of a row (only \texttt{valid} toggles), and Path-2 reads it directly.
        This removes one register bank.
  \item \textbf{Path-2 datapath left functionally unchanged} to limit the blast
        radius of the change.
\end{itemize}

\paragraph{Files changed.}
\texttt{RTL/PPU/PPU.v} (pipeline), \texttt{RTL/TOP/top.v} (stable-data feeder),
\texttt{testbench/tb\_MAC\_PPU\_hw\_flow.v} (drain + hold-data for the $+2$-cycle
latency). \texttt{tb\_top.v} needed no change (handshake-based).

\section{Verification (RTL)}

Both functional regressions pass at RTL:
\begin{itemize}
  \item \texttt{tb\_top} --- handshake-based, latency-tolerant golden check: \textbf{PASS}.
  \item \texttt{tb\_MAC\_PPU\_hw\_flow} --- rigid two-cycle protocol with the
        latency/hold-data fixes: \textbf{PASS}.
\end{itemize}

\section{Post-Fix Synthesis Result}
\label{sec:after}

% =====================================================================
% TODO: fill in after re-running synthesis on the 90nm library.
% =====================================================================

\begin{center}
\begin{tabular}{lr}
\toprule
Quantity & Value (ns) \\
\midrule
Clock period (target) & \texttt{<fill>} \\
Data required time    & \texttt{<fill>} \\
Data arrival time     & \texttt{<fill>} \\
\midrule
\textbf{Slack}        & \texttt{<fill>} \\
\bottomrule
\end{tabular}
\end{center}

\noindent New critical path: \texttt{<fill: start/end point>}.

\paragraph{Expected next bottleneck.} Once Path-1 closes, the new critical path
is most likely in Path-2: the \texttt{LPW\_reciprocal} followed by the constant
divider \texttt{path2\_scale\_numer / 127}. If it violates, the same pipelining
approach applies.

\end{document}
